\documentclass[10pt,aspectratio=169]{beamer}
\usepackage{bma01}

\title[Curvas paramétricas]{Curvas paramétricas}
\subtitle{Ecuaciones paramétricas y las derivadas\\[2pt]\small Semana 13 -- Sesión 2}

\begin{document}

\begin{frame}
  \titlepage
\end{frame}

\begin{frame}{Objetivos de la sesión}
  \begin{enumerate}
    \item Describir curvas mediante \resaltar{ecuaciones paramétricas}; reconocer su orientación.
    \item \resaltar{Eliminar el parámetro} para obtener la ecuación cartesiana.
    \item Calcular $\dfrac{dy}{dx}$ y $\dfrac{d^2y}{dx^2}$ para curvas paramétricas.
    \item Hallar rectas tangentes, tangentes horizontales y verticales.
  \end{enumerate}
  \begin{block}{Plan (2 horas)}
    \begin{tabular}{ll}
      0:00 -- 0:35 & Ecuaciones paramétricas: ejemplos, orientación, eliminación del parámetro\\
      0:35 -- 1:10 & Derivada $\frac{dy}{dx}$; recta tangente\\
      1:10 -- 1:30 & Tangentes horizontales y verticales\\
      1:30 -- 1:45 & Segunda derivada y concavidad\\
      1:45 -- 2:00 & Ejercicios y resumen
    \end{tabular}
  \end{block}
\end{frame}

\section{Ecuaciones paramétricas}

\begin{frame}{Curvas paramétricas}
  \begin{columns}
    \begin{column}{0.55\textwidth}
      \begin{definicion}
        Una \resaltar{curva paramétrica} es el conjunto de puntos
        \[ C=\{(x(t),y(t)):\ t\in I\}, \]
        donde $x=f(t)$, $y=g(t)$ son funciones del \resaltar{parámetro} $t$. El sentido en que crece $t$ da la \resaltar{orientación}.
      \end{definicion}
      \begin{ejemplo}[1: trayectoria de un proyectil]
        $x=v_0\cos\alpha\,t$, \ $y=v_0\sen\alpha\,t-\frac g2t^2$.\\
        $t$ es el tiempo: la curva dice \textbf{dónde} y \textbf{cuándo}.
      \end{ejemplo}
      \small Ventaja: describen curvas que no son gráficas de funciones (circunferencias, lazos) y movimientos.
    \end{column}
    \begin{column}{0.45\textwidth}
      \centering
      \begin{tikzpicture}
        \begin{axis}[bma, xmin=0, xmax=4.8, ymin=0, ymax=2.3, width=6.5cm, height=4.6cm, xtick=\empty, ytick=\empty, samples=60, domain=0:1.237]
          \addplot[curva, -{Stealth[length=6pt]}] ({7*x*0.5},{7*x*0.866-4.9*x^2});
          \addplot[only marks, mark=*, uniGranate, samples at={0,0.3,0.6,0.9,1.2}] ({7*x*0.5},{7*x*0.866-4.9*x^2});
          \node[uniGranate] at (axis cs:0.6,1.6) {\scriptsize $t=0.3$};
        \end{axis}
      \end{tikzpicture}

      \small Puntos cada $0.3$ s: la partícula frena al subir y acelera al bajar.
    \end{column}
  \end{columns}
\end{frame}

\begin{frame}{Ejemplos clásicos}
  \begin{columns}
    \begin{column}{0.33\textwidth}
      \centering
      \begin{tikzpicture}
        \begin{axis}[bma, xmin=-3.3, xmax=3.3, ymin=-2.3, ymax=2.3, axis equal image, width=5cm, height=4.2cm, xtick={3}, ytick={2}, samples=100]
          \addplot[curva, domain=0:340, -{Stealth[length=5pt]}] ({3*cos(x)},{2*sin(x)});
        \end{axis}
      \end{tikzpicture}\\
      \small \textbf{Elipse}: $x=3\cos t$, $y=2\sen t$. \\ $\frac{x^2}9+\frac{y^2}4=1$.
    \end{column}
    \begin{column}{0.33\textwidth}
      \centering
      \begin{tikzpicture}
        \begin{axis}[bma, xmin=0, xmax=13, ymin=0, ymax=2.6, axis equal image, width=5.4cm, height=3.5cm, xtick={6.2832,12.566}, xticklabels={$2\pi$,$4\pi$}, ytick={2}, samples=200]
          \addplot[curva2, domain=0:720] ({rad(x)-sin(x)},{1-cos(x)});
        \end{axis}
      \end{tikzpicture}\\
      \small \textbf{Cicloide}: $x=t-\sen t$, $y=1-\cos t$ (punto de una rueda).
    \end{column}
    \begin{column}{0.33\textwidth}
      \centering
      \begin{tikzpicture}
        \begin{axis}[bma, xmin=-1.3, xmax=1.3, ymin=-1.3, ymax=1.3, axis equal image, width=4.4cm, height=4.4cm, xtick=\empty, ytick=\empty, samples=400]
          \addplot[uniVerde, very thick, domain=0:360] ({sin(3*x)},{sin(2*x)});
        \end{axis}
      \end{tikzpicture}\\
      \small \textbf{Lissajous}: $x=\sen3t$, $y=\sen2t$ (osciloscopio).
    \end{column}
  \end{columns}
  \begin{ejemplo}[2: misma curva, distinta parametrización]
    $(\cos t,\sen t)$, $t\in[0,2\pi]$ y $(\cos2t,\sen2t)$, $t\in[0,\pi]$ describen la misma circunferencia, pero la segunda se recorre el doble de rápido.
    \ Con $(\sen t,\cos t)$ se recorre en sentido \textbf{horario}.
  \end{ejemplo}
\end{frame}

\begin{frame}{Eliminación del parámetro}
  \begin{ejemplo}[3]
    $x=t^2-2t$, \ $y=t+1$: \ $t=y-1$, \ $x=(y-1)^2-2(y-1)=y^2-4y+3$. Parábola horizontal.
  \end{ejemplo}
  \begin{ejemplo}[4: identidades trigonométricas]
    $x=1+2\cos t$, \ $y=-1+2\sen t$: \ $(x-1)^2+(y+1)^2=4$. Circunferencia de centro $(1,-1)$ y radio $2$.
  \end{ejemplo}
  \begin{ejemplo}[5: proyectil]
    $t=\dfrac{x}{v_0\cos\alpha}$ \ $\Rightarrow$ \ $y=x\tan\alpha-\dfrac{g}{2v_0^2\cos^2\alpha}x^2$: \ una parábola.
  \end{ejemplo}
  \begin{observacion}
    Al eliminar el parámetro se pierde la información temporal y, a veces, la porción de curva recorrida:
    $x=\sen^2t$, $y=\cos^2t$ da el \emph{segmento} $x+y=1$, $0\le x\le1$, no toda la recta.
  \end{observacion}
\end{frame}

\section{Derivadas}

\begin{frame}{La derivada $\frac{dy}{dx}$}
  \begin{teorema}
    Si $x=f(t)$, $y=g(t)$ son derivables y $f'(t)\neq0$:
    \[ \frac{dy}{dx}=\frac{dy/dt}{dx/dt}=\frac{g'(t)}{f'(t)}. \]
  \end{teorema}
  \textbf{Justificación.} Si cerca de $t_0$ podemos escribir $y=F(x)$, por la cadena $\frac{dy}{dt}=F'(x)\frac{dx}{dt}$; despejamos $F'(x)$.
  (Si $f'(t_0)\ne0$, $f$ es monótona cerca de $t_0$ y tiene inversa: así existe tal $F$.)
  \begin{ejemplo}[6]
    $x=t^2$, $y=t^3-3t$: \ $\dfrac{dy}{dx}=\dfrac{3t^2-3}{2t}$. \ En $t=2$: punto $(4,2)$, pendiente $\frac{9}{4}$.
    \ Tangente: $y-2=\frac94(x-4)$.
  \end{ejemplo}
\end{frame}

\begin{frame}{Interpretación: el vector velocidad}
  \begin{columns}
    \begin{column}{0.55\textwidth}
      Si $(x(t),y(t))$ es la posición de una partícula, el \resaltar{vector velocidad} es
      \[ \mathbf v(t)=\big(x'(t),\,y'(t)\big), \]
      tangente a la trayectoria. Su pendiente es $\frac{y'}{x'}=\frac{dy}{dx}$ y su \resaltar{rapidez} es $\sqrt{x'^2+y'^2}$.
      \begin{ejemplo}[7]
        Circunferencia $x=R\cos\omega t$, $y=R\sen\omega t$:
        $\mathbf v=(-R\omega\sen\omega t,\ R\omega\cos\omega t)$, rapidez $R\omega$ (constante), perpendicular al radio.
        \ $\frac{dy}{dx}=-\ctg\omega t=-\frac xy$.
      \end{ejemplo}
    \end{column}
    \begin{column}{0.45\textwidth}
      \centering
      \begin{tikzpicture}[scale=1.3]
        \draw[thick, uniAzul] (0,0) circle (1.3);
        \draw[->] (-1.7,0) -- (1.8,0); \draw[->] (0,-1.7) -- (0,1.8);
        \coordinate (P) at (40:1.3);
        \draw (0,0) -- (P);
        \draw[->, very thick, uniGranate] (P) -- ++(130:0.9) node[above] {\scriptsize $\mathbf v$};
        \fill (P) circle (1.5pt);
        \node at (20:0.55) {\scriptsize $\omega t$};
      \end{tikzpicture}
    \end{column}
  \end{columns}
\end{frame}

\section{Tangentes horizontales y verticales}

\begin{frame}{Tangentes horizontales y verticales}
  \small
  \begin{block}{Criterio}
    \begin{itemize}
      \item \textbf{Horizontal}: $\dfrac{dy}{dt}=0$ y $\dfrac{dx}{dt}\neq0$.
      \item \textbf{Vertical}: $\dfrac{dx}{dt}=0$ y $\dfrac{dy}{dt}\neq0$.
      \item Si ambas se anulan, analizar con límites (puede haber una \textbf{cúspide}).
    \end{itemize}
  \end{block}
  \begin{columns}
    \begin{column}{0.55\textwidth}
      \begin{ejemplo}[8]
        $x=t^2$, $y=t^3-3t$:
        \begin{itemize}
          \item Horizontal: $3t^2-3=0$, $t=\pm1$: puntos $(1,\mp2)$.
          \item Vertical: $2t=0$, $t=0$: punto $(0,0)$.
          \item La curva pasa \textbf{dos veces} por $(3,0)$ ($t=\pm\sqrt3$), con pendientes $\pm\sqrt3$: dos tangentes distintas.
        \end{itemize}
      \end{ejemplo}
    \end{column}
    \begin{column}{0.45\textwidth}
      \centering
      \begin{tikzpicture}
        \begin{axis}[bma, xmin=-0.5, xmax=4.5, ymin=-3.2, ymax=3.2, width=5.5cm, height=4.8cm, xtick={1,3}, ytick={-2,2}, samples=150]
          \addplot[curva, domain=-2.1:2.1] ({x^2},{x^3-3*x});
          \addplot[curva2, thin, domain=0:2] {2};
          \addplot[curva2, thin, domain=0:2] {-2};
          \addplot[uniVerde, thin, domain=1.5:4.5] {1.732*(x-3)};
          \addplot[uniVerde, thin, domain=1.5:4.5] {-1.732*(x-3)};
          \addplot[only marks, mark=*, uniGranate] coordinates {(1,2) (1,-2) (0,0) (3,0)};
        \end{axis}
      \end{tikzpicture}
    \end{column}
  \end{columns}
\end{frame}

\section{Segunda derivada}

\begin{frame}{La segunda derivada}
  \begin{teorema}
    \[ \frac{d^2y}{dx^2}=\frac{d}{dx}\left(\frac{dy}{dx}\right)=\frac{\dfrac{d}{dt}\left(\dfrac{dy}{dx}\right)}{\dfrac{dx}{dt}}. \]
  \end{teorema}
  \alert{¡Cuidado!} \ $\dfrac{d^2y}{dx^2}\neq\dfrac{d^2y/dt^2}{d^2x/dt^2}$.
  \begin{ejemplo}[9]
    $x=t^2$, $y=t^3-3t$: \ $\dfrac{dy}{dx}=\dfrac32t-\dfrac{3}{2t}$, \ $\dfrac{d}{dt}\left(\dfrac{dy}{dx}\right)=\dfrac32+\dfrac{3}{2t^2}$.
    \[ \frac{d^2y}{dx^2}=\frac{\frac32+\frac{3}{2t^2}}{2t}=\frac{3(t^2+1)}{4t^3}. \]
    Cóncava hacia arriba si $t>0$ (rama inferior derecha recorrida hacia arriba) y hacia abajo si $t<0$.
  \end{ejemplo}
\end{frame}

\section{Ejercicios}

\begin{frame}{Ejercicios propuestos}
  \small
  \begin{ejercicio}
    \begin{enumerate}
      \item Esboce, indique la orientación y elimine el parámetro: (a) $x=2t-1$, $y=t^2$; \ (b) $x=e^t$, $y=e^{-2t}$; \ (c) $x=\sec t$, $y=\tan t$;
        \ (d) $x=\cos^3t$, $y=\sen^3t$.
      \item Halle la recta tangente: (a) $x=t^3+1$, $y=t^2-t$ en $t=1$; \ (b) $x=e^t\cos t$, $y=e^t\sen t$ en $t=0$.
      \item Halle los puntos con tangente horizontal y vertical de: (a) $x=t^3-3t$, $y=t^2$; \ (b) $x=\cos t$, $y=\sen2t$.
      \item Halle $\frac{d^2y}{dx^2}$ y estudie la concavidad: (a) $x=t-t^2$, $y=t-t^3$; \ (b) $x=2\cos t$, $y=3\sen t$.
      \item Un proyectil sale con $v_0=30$ m/s y $\alpha=60^\circ$ ($g=9.8$). Halle la altura máxima, el alcance y el ángulo de impacto.
      \item Muestre que $x=a\cosh t$, $y=b\senh t$ parametriza una rama de la hipérbola $\frac{x^2}{a^2}-\frac{y^2}{b^2}=1$.
      \item Pruebe que la tangente a la cicloide en cada punto pasa por el punto más alto de la rueda en ese instante.
    \end{enumerate}
  \end{ejercicio}
\end{frame}

\begin{frame}{Resumen}
  \begin{itemize}
    \item Curva paramétrica: $x=f(t)$, $y=g(t)$; orientación según $t$. Eliminar $t$ da la ecuación cartesiana (con cuidado).
    \item $\dfrac{dy}{dx}=\dfrac{y'(t)}{x'(t)}$; \ $\mathbf v=(x',y')$ es tangente a la curva.
    \item Horizontal: $y'=0$, $x'\neq0$. \ Vertical: $x'=0$, $y'\neq0$. \ Ambas nulas: analizar.
    \item $\dfrac{d^2y}{dx^2}=\dfrac{\frac{d}{dt}(dy/dx)}{dx/dt}$.
  \end{itemize}
  \vspace{3mm}
  \textbf{Próxima sesión:} más ejemplos de curvas paramétricas (cicloide, cúspides, lazos, Lissajous).
\end{frame}

\end{document}
